Un controlador conecta una orden de rampa en \(t=2\); antes de ese instante la orden es cero. El tiempo y la orden están normalizados y son adimensionales. En el programa A, al conectarse se usa la lectura del reloj que funciona desde \(t=0\). En el programa B se usa un reloj que se pone a cero al conectarse. La orden es igual a la lectura del reloj correspondiente.

Construya las dos señales, calcule sus transformadas de Laplace e indique para qué \(s\) reales existen. Al aplicar un desplazamiento, identifique qué función del tiempo transcurrido está trasladando. Explique la diferencia entre las órdenes recibidas por el actuador justo después de conectar. ¿Impiden una transformada el salto finito o la elección del valor de la señal en el instante exacto de conexión?
